%% ANT LABS Technical Documentation		
%% Software Manual and Technical Document Template	
%% 									
%% This provides an example of a software manual created in Overleaf.

\documentclass{ol-softwaremanual}
%---- Required Packages and Functions ----

\usepackage{graphicx}  % For including images
\usepackage{fancyhdr}  % For custom headers and footers
\usepackage{microtype} % For typographical enhancements
\usepackage{minted}    % For code listings
\usepackage{amsmath}   % For equations and mathematics
\usepackage{hyperref}  % For hyperlinks
\usepackage[utf8]{inputenc}
\usepackage{graphicx}  % For including images
\usepackage{fancyhdr}  % For custom headers and footers
\usepackage{microtype} % For typographical enhancements
\usepackage{minted}    % For code listings
\usepackage{amsmath}   % For equations and mathematics
\usepackage{hyperref}  % For hyperlinks
\usepackage{array}           % For table customization
\usepackage{booktabs}        % For better table lines

% Set page size and margins
\usepackage[a4paper,top=3cm,bottom=3cm,left=2cm,right=2cm,headheight=50pt,headsep=20pt]{geometry}

% Custom header
\pagestyle{fancy}
\fancyhf{}  % Clear all header and footer fields
\fancyhead[L]{%
  \raisebox{-\height}{\includegraphics[width=3.5cm]{ANTLABs_logo.png}}%
   \textit{  Document Name}
}

% Custom footer
\fancyfoot[L]{%
  \parbox{15cm}{%
    \centering \thepage\\
    \scriptsize
    \textit{Confidentiality Notice: This document contains confidential and proprietary information of ANTLABS. Unauthorized use, disclosure, or distribution of this document or its contents is strictly prohibited. If you are not the intended recipient, please notify ANTLABS immediately and delete all copies of this document.}
  }
}

% Custom macros used in this example document
\newcommand{\doclink}[2]{\href{#1}{#2}\footnote{\url{#1}}}
\newcommand{\cs}[1]{\texttt{\textbackslash #1}}


% Frontmatter data; appears on title page
\title{Technical Document \\Template}
\version{1.0.0}
\author{ANTLabs}
% Define a page style with logo
\softwarelogo{\includegraphics[width=8cm]{ANTLABs_logo.png}}

\begin{document}

\maketitle
{ \centering  
    {\huge \bfseries Project Name \par}
    \vspace{2cm}
    
    \begin{tabular}{>{\bfseries}l l}
    Client Name & Altius NXT \\
    Project Name / Engagement ID & Data Extraction and Mapping Using AI \\
    Version & 1.0 \\
    \end{tabular}
}  
    \vfill
\section*{Revision History}
\noindent
\begin{tabular}{|>{\bfseries}l|l|l|l|l|l|}
    \hline
    Date & Version & Description & Author & Authorized By & Approved By \\
    \hline
    14/03/2024 & 1.0 & Draft version & Name & & \\
    \hline
    26/05/2024 & 1.1 & Edit option & Name & & \\
    \hline
\end{tabular}
    \vfill
    
 \section*{Revision Details}
\begin{tabular}{>{\bfseries}l l}
    Version & Details \\
    \midrule
    1.1 & Edit option changed to spreadsheet \\
    \end{tabular}
 
    \vfill
\newpage

\tableofcontents
\listoflistings
\newpage

\section{Introduction}

This document is an example of a software manual created using \LaTeX\ and Overleaf. It provides a nice clean format for producing this type of technical document. 

Simply start writing your document and use the Recompile button to view the updated PDF preview. Examples of commonly used commands and features can be found in the sections below, to help you get started.

To change the logo shown on the front page, simply upload your own logo image file (via the file-tree menu) to replace the default logo.pdf file. 

Once you're familiar with the editor, you can find various project settings in the Overleaf menu, accessed via the button in the very top left of the editor. To view tutorials, user guides, and further documentation, please visit our \doclink{https://www.overleaf.com/learn}{help library}. If you haven't used \LaTeX\ before, our \doclink{https://www.overleaf.com/learn/latex/Learn_LaTeX_in_30_minutes}{Learn \LaTeX\ in 30 minutes} tutorial is an excellent place to start.

For large documents, it can be useful to split the document into multiple source files.
Double-clicking on text in Overleaf's PDF preview will take you directly to the corresponding code in the editor.

\section{Including code samples}

\input{code-samples}

\section{Some further examples to get started}

\input{further-examples}

\end{document}
